% Appendix: prompts and generated material for the complex tasks.
% Included from proposal-minimal.tex.

\section{Original rule prompts}\label{app:orig}
The rules-only prompts of \citet{fu2026rules}, from their released code (\url{https://github.com/tinlaboratory/rules-vs-examples}, \texttt{rules-mode-prompts/}, MIT licence). Text in curly braces is filled per item. Tapatan uses the final-board-state prompt, whose levels differ only in its parameters: board size 3, 5, 7 and line length 3, 4, 5 for Easy, Medium, Hard \citep[Table~1]{fu2026rules}; the version in \citet[App.~F]{fu2026rules} also states the maximum number of pieces per player and that at most one player has a winning line. The Medium and Hard prompts of Lexical Category Inference are identical within each logic variant.

\subsection*{Set Game}

\paragraph{Easy.}
\begin{promptbox}
You will be shown 9 cards on the board and have to select 3 cards that form a GAME-SET based on the following rules.

RULES:
1. Each card has two attributes: an animal and a biome.
2. There are 10 types of animals: armadillo, axolotl, bat, cat, fox, horse, rabbit, sheep, turtle, wolf.
3. There are 10 types of biomes: plain, savanna, desert, swamp, forest, jungle, taiga, hill, badland, tundra.
4. A GAME-SET is a set of three cards: For each attribute (animal, biome), the three cards must be either ALL the SAME or ALL DIFFERENT. e.g. if 2 of the cards have the same value, and 1 of them has a different value, the set is NOT valid.

Here is the board:

{board}

Based on the rules, tell me which three cards here constitute a GAME-SET. Type the cards in the following format:
 \boxed{First card: CARD1
Second card: CARD2
Third card: CARD3}
\end{promptbox}

\paragraph{Medium.}
\begin{promptbox}
You will be shown 9 cards on the board and have to select 3 cards that form a GAME-SET based on the following rules.

RULES:
1. Each card has three attributes: an animal, a biome, and a food.
2. There are 10 types of animals: armadillo, axolotl, bat, cat, fox, horse, rabbit, sheep, turtle, wolf.
3. There are 10 types of biomes: plain, savanna, desert, swamp, forest, jungle, taiga, hill, badland, tundra.
4. There are 10 types of food: apple, potato, bread, cake, carrot, mutton, beef, cookie, pie, melon.
5. A GAME-SET is a set of three cards: For each attribute (animal, biome, food), the three cards must be either ALL the SAME or ALL DIFFERENT. e.g. if 2 of the cards have the same value, and 1 of them has a different value, the set is NOT valid.

Here is the board:

{board}

Based on the rules, tell me which three cards here constitute a GAME-SET. Type the cards in the following format:
 \boxed{First card: CARD1
Second card: CARD2
Third card: CARD3}
\end{promptbox}

\paragraph{Hard.}
\begin{promptbox}
You will be shown 9 cards on the board and have to select 3 cards that form a GAME-SET based on the following rules.

RULES:
1. Each card has a number and three attributes: an animal, a biome, and a food.
2. Numbers range from 1 to 12.
3. There are 10 types of animals: armadillo, axolotl, bat, cat, fox, horse, rabbit, sheep, turtle, wolf.
4. There are 10 types of biomes: plain, savanna, desert, swamp, forest, jungle, taiga, hill, badland, tundra.
5. There are 10 types of food: apple, potato, bread, cake, carrot, mutton, beef, cookie, pie, melon.
6. A GAME-SET is a set of three cards: For each attribute (animal, biome, food), the three cards must be either ALL the SAME or ALL DIFFERENT. e.g. if 2 of the cards have the same value, and 1 of them has a different value, the set is NOT valid.
7. But only for the number, 2 of the cards should have the same number, and 1 of them should have a different number in order for the set to be valid.

Here is the board:

{board}

Based on the rules, tell me which three cards here constitute a GAME-SET. Type the cards in the following format:
 \boxed{First card: CARD1
Second card: CARD2
Third card: CARD3}
\end{promptbox}

\subsection*{Tapatan}

\paragraph{all levels.}
\begin{promptbox}
You will be learning a new board game. You will then be given one final board state from that game. Determine the current game outcome.

RULES:
1. The board is a {board_size}x{board_size} grid of points with coordinates (x,y), where 0 <= x,y < {board_size}.
2. A board cell contains A for player A's piece, B for player B's piece, or . for an empty point.
3. Orthogonal neighbors are always adjacent. Diagonal neighbors are adjacent only from even-parity points where (x+y) is even: (x,y) connects to (x±1,y±1) when in bounds.
4. A player wins if they have {line_length} pieces in a straight line (horizontal, vertical, or diagonal) and each consecutive pair of points in that line is adjacent under rule 3.
5. If A has a valid winning line, output A win. If B has a valid winning line, output B win. If neither player has a valid winning line, output continue.
6. You must output a label that is exactly one of: A win, B win, or continue.

BOARD:
{board_state}

LABEL:
\end{promptbox}

\subsection*{Operator Function}

\paragraph{Easy.}
\begin{promptbox}
You will be shown a problem that uses a newly defined operator function based on the following rules and need to calculate the answer for the problem.

RULES:
1. The new operator function f is composed of two arithmetic operators, and its inputs are integers.
2. Let x, y, z be inputs of the function.
3. To compute f(x,y,z), first multiply x and y.
4. Then, add z to the result from the previous operation.
Problem: \({problem}\)
Answer (place the result in \boxed{}):\boxed{}
\end{promptbox}

\paragraph{Medium.}
\begin{promptbox}
You will be shown a problem that uses a newly defined operator function based on the following rules and need to calculate the answer for the problem.

RULES:
1. The new operator function g is composed of three arithmetic operators, and its inputs are integers.
2. Let x, y, z, a be inputs of the function.
3. To compute g(x,y,z,a), first subtract y from x.
4. Second, multiply the result from the previous operation by z.
5. Last, adds a to the result of the previous operation.
Problem: \({problem}\)
Answer (place the result in \boxed{}):\boxed{}
\end{promptbox}

\paragraph{Hard.}
\begin{promptbox}
You will be shown a problem that uses a newly defined operator function based on the following rules and need to calculate the answer for the problem.

RULES:
1. The new operator function h is composed of four arithmetic operators, and its inputs are integers.
2. Let x, y, z, a, b be inputs of the function.
3. To compute h(x,y,z,a,b), first add x and y.
4. Second, multiply the result from the previous operation by z.
5. Third, subtract a from the result of the previous operation.
6. Last, take the remainder when the result is divided by b.
Problem: \({problem}\)
Answer (place the result in \boxed{}):\boxed{}
\end{promptbox}

\subsection*{Noun Class Agreement}

\paragraph{Easy.}
\begin{promptbox}
You will be given one sentence in a synthetic language. Decide whether the sentence is grammatical.

RULES:
1. There are 2 noun classes: A and B.
2. Each noun phrase has the form DET NOUN.
3. The class A determiner is ka, and the class B determiner is ti.
4. The class A nouns are tac, sular, fep, wug. The class B nouns are bim, noko, glarn, zesh.
5. Sentences follow the template DET NOUN VERB DET NOUN.
6. Verbs are plain roots like dax, miv, lorp and never change.
7. A sentence is grammatical if and only if, in each noun phrase, the determiner matches the noun's class.
8. Your answer must be exactly one word: Yes or No.

Sentence: {sentence}
Answer:
\end{promptbox}

\paragraph{Medium.}
\begin{promptbox}
You will be given one sentence in a synthetic language. Decide whether the sentence is grammatical.

RULES:
1. There are 4 noun classes: A, B, C, D.
2. A noun phrase has the form DET ADJ-SUFF NOUN.
3. The class A determiner is ka, the class B determiner is ti, the class C determiner is su, and the class D determiner is vo.
4. Adjectives have the form STEM-SUFFIX. The suffix must match the noun's class.
5. The class A adjective suffix is -en, the class B adjective suffix is -os, the class C adjective suffix is -im, and the class D adjective suffix is -at.
6. Example adjective tokens are glim-en, prun-os, fen-im, zay-at.
7. The class A nouns are tav, fap. The class B nouns are bim, glarn. The class C nouns are sular, wug. The class D nouns are noko, zesh.
8. Sentences follow the template DET ADJ-SUFF NOUN VERB DET ADJ-SUFF NOUN.
9. Verbs are plain roots like dax, miv, lorp and never change.
10. A sentence is grammatical if and only if, in each noun phrase, both the determiner and the adjective suffix match the noun's class.
11. Your answer must be exactly one word: Yes or No.

Sentence: {sentence}
Answer:
\end{promptbox}

\paragraph{Hard.}
\begin{promptbox}
You will be given one sentence in a synthetic language. Decide whether the sentence is grammatical.

RULES:
1. There are 6 noun classes: A, B, C, D, E, F.
2. A noun phrase has the form DET ADJ-SUFF NOUN.
3. The class A determiner is ka, the class B determiner is ti, the class C determiner is su, the class D determiner is vo, the class E determiner is ne, and the class F determiner is la.
4. Adjectives have the form STEM-SUFFIX. The suffix must match the noun's class.
5. The class A adjective suffix is -en, the class B adjective suffix is -os, the class C adjective suffix is -im, the class D adjective suffix is -at, the class E adjective suffix is -uk, and the class F adjective suffix is -esh.
6. The verb has the form PREFIX-ROOT-SUFFIX.
7. The prefix must match the class of the subject noun, and the suffix must match the class of the object noun.
8. The class A subject prefix is ge-, the class B subject prefix is du-, the class C subject prefix is ri-, the class D subject prefix is zo-, the class E subject prefix is pa-, and the class F subject prefix is li-.
9. The class A object suffix is -an, the class B object suffix is -eb, the class C object suffix is -ig, the class D object suffix is -ot, the class E object suffix is -ul, and the class F object suffix is -er.
10. Example verb tokens are ge-dax-eb, pa-miv-ig, du-lorp-er.
11. The class A nouns are tav, fap. The class B nouns are bim, glarn. The class C nouns are sular, wug. The class D nouns are noko, zesh. The class E nouns are tac, fep. The class F nouns are lurn, prax.
12. Sentences follow the template DET ADJ-SUFF NOUN PREFIX-ROOT-SUFFIX DET ADJ-SUFF NOUN.
13. Verb roots are plain roots like dax, miv, lorp. Only the verb's prefix and suffix change.
14. A sentence is grammatical if and only if the subject noun phrase's determiner and adjective suffix match the subject noun's class, the object noun phrase's determiner and adjective suffix match the object noun's class, the verb prefix matches the subject noun's class, and the verb suffix matches the object noun's class.
15. Your answer must be exactly one word: Yes or No.

Sentence: {sentence}
Answer:
\end{promptbox}

\subsection*{Lexical Category Inference}

\paragraph{Easy.}
\begin{promptbox}
You will be shown a list of four words. Decide whether this list is correct based on the following rules.

RULES:
1. There are exactly four categories: {prompt_labels}.
2. Each category has the meaning shown in the category definitions below.
3. A correct list contains exactly four comma-separated items.
4. In a correct list, the first item is a member of Category 1, the second item is a member of Category 2, and so on.
5. If a category says "at least one", an item may satisfy any listed description. If it says "all", an item must satisfy every listed description.
6. Your answer must be exactly one word: Yes or No.

CATEGORY DEFINITIONS:
{label_block}

List: {candidate_block}
Answer:
\end{promptbox}

\paragraph{Medium and Hard, OR.}
\begin{promptbox}
You will be shown a list of four words. Decide whether this list is correct based on the following rules.

RULES:
1. There are exactly four categories: {prompt_labels}.
2. Each category has the meaning shown in the category definitions below.
3. A correct list contains exactly four comma-separated items.
4. In a correct list, the first item is a member of Category 1, the second item is a member of Category 2, and so on.
5. If a category lists multiple descriptions, an item is a member of that category if it satisfies at least one listed description.
6. Your answer must be exactly one word: Yes or No.

CATEGORY DEFINITIONS:
{label_block}

List: {candidate_block}
Answer:
\end{promptbox}

\paragraph{Medium and Hard, AND.}
\begin{promptbox}
You will be shown a list of four words. Decide whether this list is correct based on the following rules.

RULES:
1. There are exactly four categories: {prompt_labels}.
2. Each category has the meaning shown in the category definitions below.
3. A correct list contains exactly four comma-separated items.
4. In a correct list, the first item is a member of Category 1, the second item is a member of Category 2, and so on.
5. If a category lists multiple descriptions, an item is a member of that category only if it satisfies all listed descriptions.
6. Your answer must be exactly one word: Yes or No.

CATEGORY DEFINITIONS:
{label_block}

List: {candidate_block}
Answer:
\end{promptbox}

\section{Generation prompts}\label{app:genprompts}
Both prompts are given to GPT-5.5 through the Azure OpenAI API; \texttt{\{rule\_prompt\}} is an original prompt of App.~\ref{app:orig}.

\subsection*{Rule wordings}
\begin{promptbox}
You will be given a prompt that describes the rules of a task for a language model. Write {N} paraphrases of it.

Requirements:
1. Each paraphrase must state exactly the same rules as the original: the same constraints, conditions, quantities, names, symbols and labels. Do not add, remove, weaken or strengthen any rule, and do not add examples, hints or explanations.
2. Change the wording and the sentence structure. Reorder rules only where their order does not matter; if a rule refers to another rule by its number, keep the reference correct.
3. Copy verbatim every placeholder in curly braces (e.g. {board}), the answer-format instruction, and every output label or format string (e.g. "Yes or No", "A win, B win, or continue", \boxed{...}).
4. Keep unchanged the final lines that introduce the input and the answer (e.g. "Sentence: {sentence}" and "Answer:").
5. The paraphrases must differ from the original and from each other.

Return a JSON list of {N} strings.

Original prompt:
<<<
{rule_prompt}
>>>
\end{promptbox}

\subsection*{Term lists}
\begin{promptbox}
You will be given the rules of a task for a language model. List words that describe the task itself: what the task asks the model to do, its operation or relation, and its key concepts. For example, for the task "give the antonym of a word": opposite, reverse, antonym, contrary.

Requirements:
1. Single words only.
2. Do not list possible answers or output labels, words that only name the items of a particular input (e.g. specific animals, nouns or numbers), or placeholders.
3. Include singular and plural forms and capitalised forms. You may include translations into other languages.
4. Return a JSON list of strings.

Task rules:
<<<
{rule_prompt}
>>>
\end{promptbox}

\section{Rule wordings}\label{app:wordings}
The six wordings of each original rules prompt of App.~\ref{app:orig} at the easy level, the level in the study, generated and checked as in Sec.~\ref{sec:adapt}. Text in curly braces is filled per item.

% Generated by scripts/appendix_tex.py from materials/complex/wordings.json.

\subsection*{Set Game, easy level}

\paragraph{Wording 1.}
\begin{promptbox}
You will see 9 cards on the board, and your task is to choose 3 cards that make a GAME-SET according to these rules.

RULES:
1. Every card has two attributes: an animal and a biome.
2. The 10 animal types are: armadillo, axolotl, bat, cat, fox, horse, rabbit, sheep, turtle, wolf.
3. The 10 biome types are: plain, savanna, desert, swamp, forest, jungle, taiga, hill, badland, tundra.
4. A GAME-SET consists of three cards: For each attribute (animal, biome), the values across the three cards must be either ALL the SAME or ALL DIFFERENT. For example, if 2 cards share the same value and 1 card has a different value, the set is NOT valid.

Here is the board:

{board}

Based on the rules, tell me which three cards here constitute a GAME-SET. Type the cards in the following format:
 \boxed{First card: CARD1
Second card: CARD2
Third card: CARD3}
\end{promptbox}

\paragraph{Wording 2.}
\begin{promptbox}
You are given a board with 9 cards and must pick 3 cards that constitute a GAME-SET under the rules below.

RULES:
1. Each card contains two attributes: an animal and a biome.
2. The animals come in 10 types: armadillo, axolotl, bat, cat, fox, horse, rabbit, sheep, turtle, wolf.
3. The biomes come in 10 types: plain, savanna, desert, swamp, forest, jungle, taiga, hill, badland, tundra.
4. A GAME-SET is a set of three cards: For each attribute (animal, biome), the three cards must have values that are either ALL the SAME or ALL DIFFERENT. For instance, when 2 cards have the same value and 1 card has another value, the set is NOT valid.

Here is the board:

{board}

Based on the rules, tell me which three cards here constitute a GAME-SET. Type the cards in the following format:
 \boxed{First card: CARD1
Second card: CARD2
Third card: CARD3}
\end{promptbox}

\paragraph{Wording 3.}
\begin{promptbox}
A board containing 9 cards will be shown, and you need to select 3 cards that form a GAME-SET by following the rules.

RULES:
1. Each card is defined by two attributes: an animal and a biome.
2. There are 10 animal types: armadillo, axolotl, bat, cat, fox, horse, rabbit, sheep, turtle, wolf.
3. There are 10 biome types: plain, savanna, desert, swamp, forest, jungle, taiga, hill, badland, tundra.
4. A GAME-SET is made up of three cards: For each attribute (animal, biome), the three cards must be either ALL the SAME or ALL DIFFERENT. If, for example, 2 of the cards have one value and 1 has a different value, the set is NOT valid.

Here is the board:

{board}

Based on the rules, tell me which three cards here constitute a GAME-SET. Type the cards in the following format:
 \boxed{First card: CARD1
Second card: CARD2
Third card: CARD3}
\end{promptbox}

\paragraph{Wording 4.}
\begin{promptbox}
You will be presented with 9 cards on the board and must identify 3 cards that create a GAME-SET according to the following rules.

RULES:
1. Each card has two attributes: an animal and a biome.
2. The 10 kinds of animals are: armadillo, axolotl, bat, cat, fox, horse, rabbit, sheep, turtle, wolf.
3. The 10 kinds of biomes are: plain, savanna, desert, swamp, forest, jungle, taiga, hill, badland, tundra.
4. A GAME-SET is a group of three cards: For each attribute (animal, biome), the three cards must be either ALL the SAME or ALL DIFFERENT. e.g. when 2 of the cards share the same value and 1 of them has a different value, the set is NOT valid.

Here is the board:

{board}

Based on the rules, tell me which three cards here constitute a GAME-SET. Type the cards in the following format:
 \boxed{First card: CARD1
Second card: CARD2
Third card: CARD3}
\end{promptbox}

\paragraph{Wording 5.}
\begin{promptbox}
From the 9 cards displayed on the board, choose 3 cards that satisfy the definition of a GAME-SET using these rules.

RULES:
1. Each card includes two attributes: an animal and a biome.
2. There are 10 types of animals: armadillo, axolotl, bat, cat, fox, horse, rabbit, sheep, turtle, wolf.
3. There are 10 types of biomes: plain, savanna, desert, swamp, forest, jungle, taiga, hill, badland, tundra.
4. A GAME-SET is a collection of three cards: For each attribute (animal, biome), the three cards must be either ALL the SAME or ALL DIFFERENT. e.g. if 2 of the cards have the same value while 1 of them has a different value, the set is NOT valid.

Here is the board:

{board}

Based on the rules, tell me which three cards here constitute a GAME-SET. Type the cards in the following format:
 \boxed{First card: CARD1
Second card: CARD2
Third card: CARD3}
\end{promptbox}

\paragraph{Wording 6.}
\begin{promptbox}
You will be shown a board of 9 cards, and you have to determine 3 cards that form a GAME-SET based on the rules below.

RULES:
1. Each card has two attributes: an animal and a biome.
2. The animal attribute has 10 possible types: armadillo, axolotl, bat, cat, fox, horse, rabbit, sheep, turtle, wolf.
3. The biome attribute has 10 possible types: plain, savanna, desert, swamp, forest, jungle, taiga, hill, badland, tundra.
4. A GAME-SET is a set containing three cards: For each attribute (animal, biome), the three cards must be either ALL the SAME or ALL DIFFERENT. e.g. if 2 of the cards have matching values and 1 of them has a different value, the set is NOT valid.

Here is the board:

{board}

Based on the rules, tell me which three cards here constitute a GAME-SET. Type the cards in the following format:
 \boxed{First card: CARD1
Second card: CARD2
Third card: CARD3}
\end{promptbox}

\subsection*{Operator Function, easy level}

\paragraph{Wording 1.}
\begin{promptbox}
A problem involving a newly defined operator function will be presented, and you must determine its answer using these rules.

RULES:
1. The new operator function f consists of two arithmetic operators, and all of its inputs are integers.
2. Let x, y, z be the inputs to the function.
3. When computing f(x,y,z), begin by multiplying x and y.
4. After that, add z to the outcome of the preceding operation.
Problem: \({problem}\)
Answer (place the result in \boxed{}):\boxed{}
\end{promptbox}

\paragraph{Wording 2.}
\begin{promptbox}
You need to solve a problem that uses a newly introduced operator function defined according to the rules below.

RULES:
1. The new operator function f is made from two arithmetic operators, and its inputs are integers.
2. Let x, y, z represent the inputs of the function.
3. To evaluate f(x,y,z), multiply x by y first.
4. Next, add z to the result obtained from the previous operation.
Problem: \({problem}\)
Answer (place the result in \boxed{}):\boxed{}
\end{promptbox}

\paragraph{Wording 3.}
\begin{promptbox}
A problem will be given that relies on a newly defined operator function; compute the answer by applying the following rules.

RULES:
1. The new operator function f is formed using two arithmetic operators, and the inputs to it are integers.
2. Let x, y, z be inputs of the function.
3. For f(x,y,z), first take the product of x and y.
4. Then add z to the result of that prior operation.
Problem: \({problem}\)
Answer (place the result in \boxed{}):\boxed{}
\end{promptbox}

\paragraph{Wording 4.}
\begin{promptbox}
You will receive a problem using a newly defined operator function, and you must calculate the result according to these rules.

RULES:
1. The new operator function f is built from two arithmetic operators, and its inputs are integers.
2. Let x, y, z denote the inputs of the function.
3. In computing f(x,y,z), the first step is to multiply x and y.
4. The next step is to add z to the result from the previous operation.
Problem: \({problem}\)
Answer (place the result in \boxed{}):\boxed{}
\end{promptbox}

\paragraph{Wording 5.}
\begin{promptbox}
The problem shown will involve a newly defined operator function, and the answer must be found using the rules that follow.

RULES:
1. The new operator function f is composed of two arithmetic operators, with integer inputs.
2. Let x, y, z be the function's inputs.
3. To calculate f(x,y,z), start by multiplying x and y.
4. Then add z to the value produced by the preceding operation.
Problem: \({problem}\)
Answer (place the result in \boxed{}):\boxed{}
\end{promptbox}

\paragraph{Wording 6.}
\begin{promptbox}
Use the rules below to calculate the answer to a problem involving a newly defined operator function.

RULES:
1. The new operator function f contains two arithmetic operators, and the inputs are integers.
2. Let x, y, z be inputs of the function.
3. To compute f(x,y,z), multiply x and y before anything else.
4. Then add z to the result from the previous operation.
Problem: \({problem}\)
Answer (place the result in \boxed{}):\boxed{}
\end{promptbox}

\subsection*{Noun Class Agreement, easy level}

\paragraph{Wording 1.}
\begin{promptbox}
You will receive a single sentence in a synthetic language. Determine whether it is grammatical.

RULES:
1. The language has 2 noun classes: A and B.
2. Every noun phrase is structured as DET NOUN.
3. ka is the determiner for class A, and ti is the determiner for class B.
4. The nouns in class A are tac, sular, fep, wug. The nouns in class B are bim, noko, glarn, zesh.
5. Sentences use the pattern DET NOUN VERB DET NOUN.
6. The verbs are plain roots like dax, miv, lorp and never change.
7. A sentence is grammatical exactly when, in every noun phrase, the determiner matches the noun's class.
8. Your answer must be exactly one word: Yes or No.

Sentence: {sentence}
Answer:
\end{promptbox}

\paragraph{Wording 2.}
\begin{promptbox}
You will be provided with one sentence written in a synthetic language. Judge whether the sentence is grammatical.

RULES:
1. There are 2 noun classes: A and B.
2. A noun phrase always has the form DET NOUN.
3. The determiner for class A is ka, and the determiner for class B is ti.
4. Class A contains the nouns tac, sular, fep, wug. Class B contains the nouns bim, noko, glarn, zesh.
5. A sentence follows the template DET NOUN VERB DET NOUN.
6. Verbs are plain roots like dax, miv, lorp and never change.
7. A sentence is grammatical if and only if the determiner in each noun phrase matches the noun's class.
8. Your answer must be exactly one word: Yes or No.

Sentence: {sentence}
Answer:
\end{promptbox}

\paragraph{Wording 3.}
\begin{promptbox}
You will see one sentence in a synthetic language. Decide if that sentence is grammatical.

RULES:
1. The noun classes are 2 in number: A and B.
2. Each noun phrase is of the form DET NOUN.
3. ka is used as the class A determiner, and ti is used as the class B determiner.
4. tac, sular, fep, wug are the class A nouns. bim, noko, glarn, zesh are the class B nouns.
5. Sentences have the template DET NOUN VERB DET NOUN.
6. Verbs are plain roots like dax, miv, lorp and never change.
7. A sentence is grammatical if and only if each noun phrase has a determiner that matches the noun's class.
8. Your answer must be exactly one word: Yes or No.

Sentence: {sentence}
Answer:
\end{promptbox}

\paragraph{Wording 4.}
\begin{promptbox}
A single sentence in a synthetic language will be given to you. Decide whether it is grammatical.

RULES:
1. There are 2 noun classes, A and B.
2. Noun phrases take the form DET NOUN.
3. The class A determiner is ka; the class B determiner is ti.
4. The class A nouns are tac, sular, fep, wug. The class B nouns are bim, noko, glarn, zesh.
5. Sentences are built according to DET NOUN VERB DET NOUN.
6. Verbs are plain roots like dax, miv, lorp and never change.
7. A sentence is grammatical if and only if, for each noun phrase, the determiner matches the noun's class.
8. Your answer must be exactly one word: Yes or No.

Sentence: {sentence}
Answer:
\end{promptbox}

\paragraph{Wording 5.}
\begin{promptbox}
You will get one sentence from a synthetic language. Determine whether the sentence is grammatical.

RULES:
1. There are 2 noun classes: A and B.
2. The form of every noun phrase is DET NOUN.
3. For class A, the determiner is ka, and for class B, the determiner is ti.
4. The class A nouns are tac, sular, fep, wug. The class B nouns are bim, noko, glarn, zesh.
5. Sentences conform to the template DET NOUN VERB DET NOUN.
6. Verbs are plain roots like dax, miv, lorp and never change.
7. A sentence is grammatical if and only if the determiner matches the noun's class in each noun phrase.
8. Your answer must be exactly one word: Yes or No.

Sentence: {sentence}
Answer:
\end{promptbox}

\paragraph{Wording 6.}
\begin{promptbox}
You will be shown one sentence in a synthetic language. Decide whether that sentence is grammatical.

RULES:
1. The language includes 2 noun classes: A and B.
2. Each noun phrase follows the form DET NOUN.
3. ka is the class A determiner, while ti is the class B determiner.
4. The class A nouns are tac, sular, fep, wug. The class B nouns are bim, noko, glarn, zesh.
5. Sentences follow this template: DET NOUN VERB DET NOUN.
6. Verbs are plain roots like dax, miv, lorp and never change.
7. A sentence is grammatical if and only if, within every noun phrase, the determiner matches the noun's class.
8. Your answer must be exactly one word: Yes or No.

Sentence: {sentence}
Answer:
\end{promptbox}

\subsection*{Tapatan, easy level}

\paragraph{Wording 1.}
\begin{promptbox}
You are going to learn a new board game. After that, you will receive a single final board state from the game. Decide the present game outcome.

RULES:
1. The board is a {board_size}x{board_size} grid of points whose coordinates are (x,y), with 0 <= x,y < {board_size}.
2. Each board cell holds A for player A's piece, B for player B's piece, or . for an empty point.
3. Orthogonal neighbors are always adjacent. Diagonal neighbors are adjacent only from even-parity points where (x+y) is even: (x,y) connects to (x±1,y±1) when in bounds.
4. A player wins when they have {line_length} pieces in one straight line (horizontal, vertical, or diagonal), with every consecutive pair of points in that line adjacent according to rule 3.
5. If A has a valid winning line, output A win. If B has a valid winning line, output B win. If neither player has a valid winning line, output continue.
6. You must output a label that is exactly one of: A win, B win, or continue.

BOARD:
{board_state}

LABEL:
\end{promptbox}

\paragraph{Wording 2.}
\begin{promptbox}
You will be taught a new board game, and then you will be shown one final board state from that game. Identify the game outcome at that moment.

RULES:
1. The board consists of a {board_size}x{board_size} grid of points at coordinates (x,y), where 0 <= x,y < {board_size}.
2. A board cell contains A for player A's piece, B for player B's piece, or . for an empty point.
3. Orthogonal neighbors are always adjacent. Diagonal neighbors are adjacent only from even-parity points where (x+y) is even: (x,y) connects to (x±1,y±1) when in bounds.
4. A player wins if they possess {line_length} pieces in a straight line (horizontal, vertical, or diagonal) and each successive pair of points in that line is adjacent under rule 3.
5. If A has a valid winning line, output A win. If B has a valid winning line, output B win. If neither player has a valid winning line, output continue.
6. You must output a label that is exactly one of: A win, B win, or continue.

BOARD:
{board_state}

LABEL:
\end{promptbox}

\paragraph{Wording 3.}
\begin{promptbox}
You will learn the rules for a new board game. You will then be provided with one final board state from it. Determine what the current outcome of the game is.

RULES:
1. The board is a grid of points with size {board_size}x{board_size} and coordinates (x,y), where 0 <= x,y < {board_size}.
2. In a board cell, A denotes player A's piece, B denotes player B's piece, and . denotes an empty point.
3. Orthogonal neighbors are always adjacent. Diagonal neighbors are adjacent only from even-parity points where (x+y) is even: (x,y) connects to (x±1,y±1) when in bounds.
4. A player wins by having {line_length} pieces in a straight line (horizontal, vertical, or diagonal), provided that each consecutive pair of points in that line is adjacent under rule 3.
5. If A has a valid winning line, output A win. If B has a valid winning line, output B win. If neither player has a valid winning line, output continue.
6. You must output a label that is exactly one of: A win, B win, or continue.

BOARD:
{board_state}

LABEL:
\end{promptbox}

\paragraph{Wording 4.}
\begin{promptbox}
You will be introduced to a new board game. Next, one final board state from the game will be given to you. Work out the current game outcome.

RULES:
1. The board is a {board_size}x{board_size} grid made of points with coordinates (x,y), where 0 <= x,y < {board_size}.
2. A board cell contains A for player A's piece, B for player B's piece, or . for an empty point.
3. Orthogonal neighbors are always adjacent. Diagonal neighbors are adjacent only from even-parity points where (x+y) is even: (x,y) connects to (x±1,y±1) when in bounds.
4. A player wins if {line_length} of their pieces form a straight line (horizontal, vertical, or diagonal) and every consecutive pair of points in that line is adjacent under rule 3.
5. If A has a valid winning line, output A win. If B has a valid winning line, output B win. If neither player has a valid winning line, output continue.
6. You must output a label that is exactly one of: A win, B win, or continue.

BOARD:
{board_state}

LABEL:
\end{promptbox}

\paragraph{Wording 5.}
\begin{promptbox}
You will receive the rules of a new board game. Then you will be given one final state of the board from that game. Determine the outcome currently shown.

RULES:
1. The board is a {board_size}x{board_size} point grid using coordinates (x,y), where 0 <= x,y < {board_size}.
2. A board cell contains A for player A's piece, B for player B's piece, or . for an empty point.
3. Orthogonal neighbors are always adjacent. Diagonal neighbors are adjacent only from even-parity points where (x+y) is even: (x,y) connects to (x±1,y±1) when in bounds.
4. A win occurs for a player if they have {line_length} pieces in a straight line (horizontal, vertical, or diagonal) and each consecutive pair of points in that line is adjacent under rule 3.
5. If A has a valid winning line, output A win. If B has a valid winning line, output B win. If neither player has a valid winning line, output continue.
6. You must output a label that is exactly one of: A win, B win, or continue.

BOARD:
{board_state}

LABEL:
\end{promptbox}

\paragraph{Wording 6.}
\begin{promptbox}
You will study a new board game and then receive one final board state from that game. Decide which current game outcome applies.

RULES:
1. The board is a {board_size}x{board_size} grid of points, each with coordinates (x,y), where 0 <= x,y < {board_size}.
2. A board cell contains A for player A's piece, B for player B's piece, or . for an empty point.
3. Orthogonal neighbors are always adjacent. Diagonal neighbors are adjacent only from even-parity points where (x+y) is even: (x,y) connects to (x±1,y±1) when in bounds.
4. A player wins if they have {line_length} pieces arranged in a straight line (horizontal, vertical, or diagonal) and each consecutive pair of points in that line is adjacent under rule 3.
5. If A has a valid winning line, output A win. If B has a valid winning line, output B win. If neither player has a valid winning line, output continue.
6. You must output a label that is exactly one of: A win, B win, or continue.

BOARD:
{board_state}

LABEL:
\end{promptbox}

\subsection*{Lexical Category Inference, easy level}

\paragraph{Wording 1.}
\begin{promptbox}
You will receive a list made up of four words. Determine whether the list is correct according to these rules.

RULES:
1. There are exactly four categories: {prompt_labels}.
2. The meaning of each category is given in the category definitions below.
3. A correct list has exactly four items separated by commas.
4. For a list to be correct, item one must belong to Category 1, item two must belong to Category 2, and this pattern continues in order.
5. When a category says "at least one", an item may match any of the listed descriptions. When it says "all", an item must match every listed description.
6. Your answer must be exactly one word: Yes or No.

CATEGORY DEFINITIONS:
{label_block}

List: {candidate_block}
Answer:
\end{promptbox}

\paragraph{Wording 2.}
\begin{promptbox}
A list of four words will be presented to you. Judge whether it is correct using the rules below.

RULES:
1. There are exactly four categories: {prompt_labels}.
2. Each category means what is specified in the category definitions below.
3. A correct list consists of exactly four comma-separated items.
4. In a correct list, the items correspond in order: the first item is in Category 1, the second item is in Category 2, and so on.
5. If a category says "at least one", an item may satisfy any listed description. If it says "all", an item must satisfy every listed description.
6. Your answer must be exactly one word: Yes or No.

CATEGORY DEFINITIONS:
{label_block}

List: {candidate_block}
Answer:
\end{promptbox}

\paragraph{Wording 3.}
\begin{promptbox}
You will be given four words as a list. Decide if that list is correct by applying the following rules.

RULES:
1. There are exactly four categories: {prompt_labels}.
2. The category definitions below state the meaning of each category.
3. To be correct, a list must include exactly four items separated by commas.
4. To be correct, the first item must be a member of Category 1, the second item must be a member of Category 2, and so forth.
5. If a category says "at least one", an item may satisfy any listed description. If it says "all", an item must satisfy every listed description.
6. Your answer must be exactly one word: Yes or No.

CATEGORY DEFINITIONS:
{label_block}

List: {candidate_block}
Answer:
\end{promptbox}

\paragraph{Wording 4.}
\begin{promptbox}
You will see a four-word list. Assess whether this list is correct under the rules that follow.

RULES:
1. There are exactly four categories: {prompt_labels}.
2. Each category has the meaning provided in the category definitions below.
3. A list is correct only if it contains exactly four comma-separated items.
4. A correct list places a Category 1 member first, a Category 2 member second, and continues in that order.
5. If a category says "at least one", an item may satisfy any listed description. If it says "all", an item must satisfy every listed description.
6. Your answer must be exactly one word: Yes or No.

CATEGORY DEFINITIONS:
{label_block}

List: {candidate_block}
Answer:
\end{promptbox}

\paragraph{Wording 5.}
\begin{promptbox}
You will be shown four words in list form. Decide whether the list qualifies as correct according to these rules.

RULES:
1. There are exactly four categories: {prompt_labels}.
2. The definitions below give the meaning for each category.
3. Exactly four comma-separated items are required for a correct list.
4. In a correct list, the first, second, and later items must respectively be members of Category 1, Category 2, and the following categories in sequence.
5. If a category says "at least one", an item may satisfy any listed description. If it says "all", an item must satisfy every listed description.
6. Your answer must be exactly one word: Yes or No.

CATEGORY DEFINITIONS:
{label_block}

List: {candidate_block}
Answer:
\end{promptbox}

\paragraph{Wording 6.}
\begin{promptbox}
A list containing four words will be provided. Determine whether it is correct based on the rules listed here.

RULES:
1. There are exactly four categories: {prompt_labels}.
2. Each category's meaning is the one shown in the category definitions below.
3. For a list to be correct, it must contain exactly four items, with commas separating them.
4. For a list to be correct, its first item must belong to Category 1, its second item must belong to Category 2, and the remaining items must follow the same order.
5. If a category says "at least one", an item may satisfy any listed description. If it says "all", an item must satisfy every listed description.
6. Your answer must be exactly one word: Yes or No.

CATEGORY DEFINITIONS:
{label_block}

List: {candidate_block}
Answer:
\end{promptbox}


\section{Term lists}\label{app:terms}
The term list of each complex task at the easy level, generated and checked as in Sec.~\ref{sec:adapt}, as words before their expansion to token ids. Tapatan's list is made from its single prompt, shared by all levels.

\begin{center}
% Generated by scripts/appendix_tex.py from materials/complex/terms.json.
\begin{footnotesize}
\begin{tabular}{@{}lp{0.7\textwidth}@{}}
\toprule
Task (easy level) & Terms \\
\midrule
Set Game & set, sets, Set, Sets, card, cards, Card, Cards, board, boards, Board, Boards, select, selects, Select, Selects, identify, identifies, Identify, Identifies, find, finds, Find, Finds, choose, chooses, Choose, Chooses, match, matches, Match, Matches, matching, Matching, valid, Valid, invalid, Invalid, attribute, attributes, Attribute, Attributes, animal, animals, Animal, Animals, biome, biomes, Biome, Biomes, same, Same, different, Different, all, All, triple, triples, Triple, Triples, triad, triads, Triad, Triads, combination, combinations, Combination, Combinations, group, groups, Group, Groups, game, games, Game, Games, gameset, gamesets, Gameset, Gamesets, rule, rules, Rule, Rules, constraint, constraints, Constraint, Constraints, value, values, Value, Values \\[2pt]
Operator Function & arithmetic, Arithmetic, operator, operators, Operator, Operators, function, functions, Function, Functions, calculation, calculations, Calculation, Calculations, compute, Compute, calculate, Calculate, evaluate, Evaluate, solve, Solve, multiplication, multiplications, Multiplication, Multiplications, multiply, Multiply, product, products, Product, Products, addition, additions, Addition, Additions, add, Add, sum, sums, Sum, Sums, integer, integers, Integer, Integers, input, inputs, Input, Inputs, result, results, Result, Results, expression, expressions, Expression, Expressions, composition, compositions, Composition, Compositions \\[2pt]
Noun Class Agreement & grammar, Grammar, grammars, Grammars, grammaticality, Grammaticality, grammatical, Grammatical, ungrammatical, Ungrammatical, syntax, Syntax, sentence, Sentence, sentences, Sentences, phrase, Phrase, phrases, Phrases, noun, Noun, nouns, Nouns, determiner, Determiner, determiners, Determiners, class, Class, classes, Classes, agreement, Agreement, match, Match, matches, Matches, matching, Matching, mismatch, Mismatch, mismatches, Mismatches, template, Template, templates, Templates, validation, Validation, classification, Classification, decision, Decision, decisions, Decisions, binary, Binary, wellformedness, Wellformedness, acceptability, Acceptability \\[2pt]
Tapatan & board, boards, Board, Boards, game, games, Game, Games, grid, grids, Grid, Grids, cell, cells, Cell, Cells, point, points, Point, Points, coordinate, coordinates, Coordinate, Coordinates, piece, pieces, Piece, Pieces, player, players, Player, Players, state, states, State, States, outcome, outcomes, Outcome, Outcomes, result, results, Result, Results, winner, winners, Winner, Winners, win, wins, Win, Wins, winning, Winning, line, lines, Line, Lines, alignment, alignments, Alignment, Alignments, sequence, sequences, Sequence, Sequences, straight, Straight, horizontal, Horizontal, vertical, Vertical, diagonal, diagonals, Diagonal, Diagonals, orthogonal, Orthogonal, neighbor, neighbors, Neighbor, Neighbors, adjacency, adjacencies, Adjacency, Adjacencies, adjacent, Adjacent, connection, connections, Connection, Connections, connectivity, Connectivity, parity, parities, Parity, Parities, even, Even, empty, Empty, validity, Validity, valid, Valid, classification, classifications, Classification, Classifications, classify, Classify, determine, Determine, detect, Detect, evaluate, Evaluate, identify, Identify, label, labels, Label, Labels \\[2pt]
Lexical Category Inference & validation, Validation, validate, Validate, verification, Verification, assessment, Assessment, decision, Decision, classification, Classification, classify, Classify, categorization, Categorization, categorize, Categorize, category, Category, categories, Categories, membership, Membership, member, Member, members, Members, list, List, lists, Lists, item, Item, items, Items, word, Word, words, Words, order, Order, ordering, Ordering, sequence, Sequence, position, Position, positions, Positions, alignment, Alignment, matching, Matching, match, Match, criteria, Criteria, criterion, Criterion, rule, Rule, rules, Rules, constraint, Constraint, constraints, Constraints, definition, Definition, definitions, Definitions, correctness, Correctness, incorrectness, Incorrectness, satisfaction, Satisfaction, validity, Validity \\[2pt]
\bottomrule
\end{tabular}
\end{footnotesize}

\end{center}
